\documentclass{article}
\usepackage{graphicx} % Required for inserting images
\usepackage{amsfonts, amsmath, amssymb, amsthm}
\usepackage{blindtext}

\newcommand{\set}[1]{\left\{#1\right\}}
\newcommand{\paren}[1]{\left (#1 \right )}
\newcommand{\abs}[1]{\left |#1 \right |}
\newcommand{\angles}[1]{\left \langle #1 \right \rangle}
\def \N {\mathbb N}
\def \R {\mathbb R}
\def \Q {\mathbb Q}
\def \Z {\mathbb Z}
\def \C {\mathbb C}


\title{Formalizing Ear Decompositions}
\author{Denny Yang}
\date{November 2026}

\begin{document}

\maketitle

\section{Introduction}
An ear decomposition, generally speaking, is a partitioning of a graph's edges into "ears" $P_0,P_1,..., P_t$. There does not seem to be a consensus on the specifics in literature; one definition would be that ears are walks with no repeated edges or vertices except possibly the endpoints and whose endpoints belong to the previous ears. Sometimes $P_0$ is specified to be a cycle.
\\ A graph is $k$-edge-connected if the graph remains connected after removing any $k-1$ edges. It is well-known that a graph is 2-edge-connected if and only if it has an ear decomposition\footnote{Wikipedia attributes this to Robbins in 1939, while Lovasz 1985, attributes this to Whitney in 1932. Chasing this down is not an important part of this writeup, so I will leave it at that}, which often works as an useful characterization in some proofs or algorithms.
\section{Lean Formalization}
Ultimately, I proved only the reverse direction of the bi-implication, that the existence of an ear decomposition(as formalized) gives that a graph is 2-edge-connected. Some work was done on the other direction, but it remains unfinished. Completed proofs (including ones intended for the forward direction but are unused) are in Finite.lean, while incomplete proofs, along with comments explaining my idea, are in FiniteUnfinished.lean.
\\I formalized partial ear decomposition as an inductive type partialEarDecomp, whose base case is a single cycle and whose inductive case is a walk+a proof that it is an ear. Each also contains the subgraph generated by the union of every ear. 
\\ This was because the Lean proof that partialEarDecomps are 2-edge-connected was based off of the inductive proof from a textbook(Digraphs: Theory), so an inductive type naturally made sense. This does limit the scope to finite graphs only, as walks in Lean are finite, so a finite inductive type cannot possibly capture every edge of a 2-edge-connected(and thus connected) infinite graph. However, the material I found all specify a finite graph. 
\\ In my original idea for this proof, I would have been showing that every vertex in a partialEarDecomp are 2-reachable within the subgraph formed by partialEarDecomp and not just 2-reachable within the whole graph; this idea quickly fell apart when I actually tried it, as edge-connectivity on subgraphs is not in Mathlib and coercing subgraphs into graphs adds much complexity. Luckily, I managed to use transitivity of edge-reachability, something which is in Mathlib, to get a far simpler proof than even what I originally intended.
\\ This proof that partialEarDecomps are 2-edge-connected naturally resulted in a proof that a ear decomposition is also 2-edge-connected.
\\ As for the reverse direction, the natural proof would be to construct such ear decomposition. The lemma nonempty\_two\_edge\_conn\_has\_cycle shows that nonempty 2-edge-connected graphs are not IsAcyclic as defined in Mathlib, from which deriving the existence of a cycle using classical logic should be simple, allowing the construction of the first ear $P_0$, or partialEarDecomp.init in my formalization
\\ From there, we can add ears with new vertices until none remain, and then add single edges between vertices until all none remain. This process must terminate since we have a finite graph and each added ear has at least one edge.
\\ The latter of the two steps is justified in natural language by the fact these single-edge walks are always paths and thus valid ears, and is formalized as single\_edge\_walks\_are\_ears
\\ The former of the two steps is justified in natural language as such: since not all vertices are yet in the ear decomposition, and the graph is connected, there must be a boundary edge between a vertex $u$ in and a vertex $w$ not in the ear decomposition. Since the graph is 2-edge-connected, we can delete this edge and still have a path between $u$ and $w$, which necessarily must not contain the edge $uw$. Trim this path starting with $u$ until we are at the vertex closest to $w$ which is in the ear decomposition. Finally, add the edge $uw$ to the trimmed path; verify it is an ear. Of this, I finished up to the trimming part, which was something too complicated for me to work with. Some sort of induction would likely work.
\section{Reflection}
Overall, I found it quite surprising how much work goes into the initial formalization. Furthermore, I felt quite restricted by the confines of Lean and Mathlib. In natural language proofs, the idea of "converting" a walk into a subgraph does not ever come up, but as they are separate types in Mathlib it requires an explicit method and a couple of membership propositions to do the same thing, and this goes doubly true for subgraph-graph conversions, which I ultimately chose to completely forgo due to the complexity of it when I found alternatives.
\\ Obviously, this projects leaves much to be completed. In addition to that, there are some natural extensions. Weak ear decompositions requires all ears other than $P_0$ to be paths, and their existence is a characterization of 2-vertex-connected graphs (with edge-case exceptions, depending on the specific definition of vertex connectivity). The natural language proofs for the two are extremely similar, so while some work will need to be done, it should be relatively simple to complete.
\section{Additional Attribution}
Much of my knowledge of the subject comes from a Spring 2026 21484 Graph Theory course. To complete this project, I used some of my own handwritten notes taken during lecture, plus whatever I remembered, as a supplement to cited literature.
\newpage
\section*{References}
Bang-Jensen, Jørgen; Gutin, Gregory (2007), "7.2 Ear Decompositions", Digraphs: Theory, Algorithms and Applications, Springer-Verlag, pp. 349–352
\\ Frank, A. Conservative weightings and ear-decompositions of graphs. Combinatorica 13, 65–81 (1993). https://doi.org/10.1007/BF01202790
\\ L. Lovasz, "Computing ears and branchings in parallel," 26th Annual Symposium on Foundations of Computer Science (sfcs 1985), Portland, OR, USA, 1985, pp. 464-467, doi: 10.1109/SFCS.1985.16.
\\Whitney, Hassler. “Non-Separable and Planar Graphs.” Transactions of the American Mathematical Society 34, no. 2 (1932): 339–62. https://doi.org/10.2307/1989545.
\end{document}
